% Transcription of folder 147, batch 03 (archivists' pages 41-60).
% Pass: Opus 5.5 (claude-opus-5-5), 2026-10-09 — first pass, unchecked against the pages by a human.
%
% All twenty pages are mathematics in his hand; none skipped. He numbers every
% other sheet (author's pp. 20-29 = archivists' pp. 41, 43, ..., 59); the
% unnumbered sheets in between continue the text. The batch opens mid-argument
% (n° 6, generalisation morphisms of graphs, formula (48) is before the batch)
% and ends mid-sentence on p. 60.
\documentclass[11pt,a4paper]{article}
\input{../preamble/grothendieck}

\folder{147}
\batch{3}
\pages{41}{60}
\foldertitle{Structure à l'infini des Mg,ν [pages 1 à 90, dont table des matières] : notes manuscrites (s.d.).}
\dating{s.d. — le groupe « Autour de La "Longue Marche" à travers la théorie de Galois » (141 à 149) est daté [à partir de 1978]-1983}
\watermark{Édition de démonstration}

\begin{document}

\note{Le lot s'ouvre au milieu d'un argument : la condition 1) de (48), à laquelle renvoie la première phrase, précède ce lot.}

\page{41}\note{p. 20 de l'auteur}
Reprenons un morphisme généralisé de graphes ordinaires $f : G' \to G$. On dit que $f$ est un \emph{morphisme de généralisation} s'il satisfait la condition 1) de (48).

Notons que $G$ et $f$ sont définis alors à isom. près par $G'$ et l'ens $\widetilde{B}' \subset \widetilde{A}'$ (stable par $\sigma_{\widetilde{A}'}$), ou ce qui revient au même, par $G'$ et la \struck{\ill{}} partie $B'$ de $A' = \widetilde{A}'/\sigma_{\widetilde{A}'}$ ($B'$ est donc un ens. d'arêtes de $G'$). En effet, soit $G'_{!}$ le sous-graphe de $G'$ défini par $S'$ et $A' \smallsetminus B' \subset A'$ (ou $\widetilde{A}' \smallsetminus \widetilde{B}' \subset \widetilde{A}'$) (i.e.\ \add{par $S'$ et} l'ens des \struck{arêtes} arcs de $G'$ qui sont au dessus des identités de $\Pi_1(G)$). Alors il résulte de la condition $S' \to S$ surjectif \add{et} les $G'(\alpha)$ connexes (\(\Leftrightarrow\) les $G'(\alpha)$ $0$-connexes) que \struck{les fibres} l'homomorphisme canonique
\[
\pi_0(G'_{!}) \longrightarrow S
\]
est bijectif, d'où
\[
(49)\qquad S \simeq \pi_0(G'_{!}),
\]
d'ailleurs on a
\[
(50)\qquad \widetilde{A} \simeq \widetilde{B}'
\]
\struck{c}compatible avec $\sigma_{\widetilde{B}'}$ et $\sigma_{\widetilde{A}}$

\page{42}
enfin $\widetilde{A} \xrightarrow{\sigma_G} S$ est déterminé par passage au quotient à partir de $\sigma_G \,|\, \widetilde{B}' : \widetilde{B}' \to S'$, par commutativité de

\begin{tikzcd}
\widetilde{B}' \arrow[r, "{\sigma_{G'} | \widetilde{B}'}"] \arrow[d, "\wr"'] & S' = \text{sommets de } G' \arrow[d] \\
\widetilde{A} \arrow[r] & S \simeq \pi_0(G'_{!})
\end{tikzcd}

\note{diagramme numéroté (51) en marge.}

Ainsi $G$ se reconstruit à partir de $G'$, et $f$ de \struck{\ill{}} bien sûr. Quant au fait que $G$ apparaît comme un quotient de $G'$, en fait c'est \uncertain{plutôt} une somme amalgamée dans la catégorie des graphes et des morphismes généralisés de graphes :

\begin{tikzcd}[column sep=large]
G'_{!} \arrow[r, hook, "\text{incl. can}"] \arrow[d, "\text{proj. can}"'] & G' \arrow[d] \\
S = \pi_0(G'_{!}) \arrow[r] & G
\end{tikzcd}

\note{diagramme numéroté (52) en marge, « cocart » inscrit au centre du carré ; sous $S = \pi_0(G'_{!})$ : « (graphe \struck{produit} \add{discret} défini par $S \overset{\text{déf}}{=} \pi_0(G'_{!})$) ».}

i.e.\ $G' \to G$ est universel pour les morphismes généralisés de \add{$G$ dans un} \uncertain{graphes} $G'$, qui transforment les arcs de $\widetilde{A}' \smallsetminus \widetilde{B}'$ en des identités. Posons donc

\page{43}\note{p. 21 de l'auteur}
\[
(53)\qquad \widetilde{C}' = \widetilde{A}' \smallsetminus \widetilde{B}', \quad C' = A' \smallsetminus B' = \widetilde{C}'/\sigma_{\widetilde{A}'},
\]
on écrira
\[
(54)\qquad G = G'/\widetilde{C}' \quad\text{ou}\quad G = G'/C'
\]
et on dira que $G$ est \emph{déduit de $G'$} en contractant les arêtes $a' \in C'$ (ou encore, les arcs $\tilde{a}' \in \widetilde{C}'$) en des pts. \uncertain{Interprétant} les morphismes généralisés de graphes, en termes des réalisations \struck{géométriques} \add{topologiques} de ceux-ci, à l'aide d'applications continues qui contractent justement certaines arêtes en des pts, on voit que cette terminologie correspond bien à une intuition topologique courante. D'ailleurs la somme amalgamée (52) existe quel que soit le sous-graphe $G'_{!}$ de $G'$ ayant mêmes sommets que $G'$, i.e.\ quel que soit l'ens d'\struck{flèches} \add{arêtes} $C' \subset A'$ de $G'_{!}$. [Il s'ensuit que tout morphisme généralisé de graphes $G' \to G$ se factorise de façon unique en un
\page{44}
composé
\[
(55)\qquad G' \to G_0 \to G,
\]
où $G' \to G_0$ est un morphisme de passage au quotient (par l'ens des arêtes de $C'$ contractées par $f$, justement) et où $G_0 \to G$ est un morphisme \add{ordinaire} de graphes.


Supposons maintenant que $G'$ soit \struck{\ill{}} \add{fini, et} provienne\add{\ill{}} d'un MD-graphe $\underline{G}'$, grâce aux données supplémentaires
\[
I' \to S', \quad S' \xrightarrow{g'} \mathbb{N},
\]
où $I'$ est un ens. fini. Alors $G$ provient canoniquement d'un MD-graphe $\underline{G}$, en posant $I = I_{\underline{G}} = I'$, et en définissant $I \to S$ comme le composé
\[
I \overset{\text{déf}}{=} I' \to S' \to S,
\]
et en définissant, pour $\forall\, \alpha \in S$
\[
(56)\qquad g_\alpha = \Bigl(\sum_{\alpha' \in S'_\alpha} g'_{\alpha'}\Bigr) + h^1(G'(\alpha)).
\]
\note{sous $S'_\alpha$ : « $=$ sommets de $G'(\alpha)$ ». Devant la ligne $I' \to S'$, un numéro biffé.}

\page{45}\note{p. 22 de l'auteur}
Il faut cependant vérifier encore la condition MD sur le graphe $\underline{G}$ obtenu. Que $\underline{G}$ soit connexe résulte aussitôt de la connexité de $\underline{G}'$. Pour $\forall\, \alpha \in S$, il faut vérifier que
\[
2g_\alpha + \underbrace{\operatorname{card} I_\alpha + \operatorname{card} \widetilde{A}_\alpha}_{\hat{\nu}_\alpha} \geqslant 3,
\]
\struck{ce qui dit que \ill{}} i.e.\ que $(g_\alpha, \hat{\nu}_\alpha)$ est \uncertain{admissible}. Or \struck{ou que $\operatorname{card} 2g_\alpha + \operatorname{card} I_\alpha \geqslant 3$}
\struck{qui dit} \add{on sait} que $(g_\alpha, \hat{\nu}_\alpha)$ est le type numérique du \uncertain{MD-graphe} $\underline{G}'(\alpha)$, dont le graphe sous-jacent est la composante connexe de $G'_{!}$ correspondant à $\alpha$, et l'ens.\ des pts « marqués » est $I'_\alpha \amalg \widetilde{B}'_\alpha$, où $I'_\alpha = I' \,|\, S'_\alpha \simeq I_\alpha$, et $\widetilde{B}'_\alpha = \widetilde{A}'_\alpha \cap \widetilde{B}'$\note{la lettre avant « $'_\alpha \cap$ » est surchargée ($B$ sur $A$ ?) ; le sens demande l'ensemble des flèches d'origine dans $S'_\alpha$.} est l'ens des flèches de $\widetilde{B}'$ d'origine dans $S'_\alpha$, l'homomorphisme de projection
\[
I(\alpha) \overset{\text{déf}}{=} I'_\alpha \amalg \widetilde{B}'_\alpha \longrightarrow S'_\alpha
\]
étant induit par $I' \to S'$ et l'application origine. Donc, en vertu de la prop.\ du n°~3, on aura $(g_\alpha, \hat{\nu}_\alpha)$ \uncertain{admissible} pourvu seulement que $\underline{G}(\alpha)$ est un MD-graphe. Mais cette condition signifie que l'on a, pour $\forall\, \beta \in S'_\alpha$

\page{46}
\[
2g'_\beta + \underbrace{\bigl(\operatorname{card} I'_\beta + \operatorname{card} \widetilde{B}'_\beta\bigr)}_{\substack{\text{poids marqué}\\ \text{de } \beta \text{ dans } \underline{G}(\alpha)}} + \underbrace{\operatorname{card} \widetilde{C}'_\beta}_{\substack{\text{ordre de}\\ \beta \text{ dans } \underline{G}(\alpha)}} \geqslant 3
\]
\note{après $I'_\beta$, quelques signes biffés.}
or cette expression s'écrit aussi, puisque $\widetilde{A}' = \widetilde{B}' \amalg \widetilde{C}'$
\[
2g'_\beta + \underbrace{\operatorname{card} I'_\beta}_{\substack{\text{poids marqué}\\ \text{de } \beta \text{ dans } \underline{G}'_\beta}} + \underbrace{\operatorname{card} \widetilde{A}'_\beta}_{\substack{\text{ordre de } \beta\\ \text{dans } \underline{G}'}} \geqslant 3,
\]
et cette condition n'est autre que la condition \add{de régularité} MD sur $\underline{G}'$.

Résumons les résultats sur les MD-graphes\note{ici et plus bas, « MD-graphe » est écrit très vite, en abrégé ; la lecture est fixée par les occurrences nettes de « MD-graphe » sur la page 51.} :

\emph{Proposition.} Soit $\underline{G}' = (G', I' \to S' \xrightarrow{g'} \mathbb{N})$, où $G' = (S', \widetilde{A}', \sigma_{\widetilde{A}'}, o_{\widetilde{A}'})$, un MD-graphe, et soit $C' \subset A' = \widetilde{A}'/\sigma_{\widetilde{A}'}$ un ens d'arêtes de $G'$. Alors on peut de façon unique \struck{\ill{}} \add{(à isom.\ unique près)} trouver un morphisme de généralisation de MD-graphes
\[
\underline{G}' \longrightarrow \underline{G}
\]

\page{47}\note{p. 23 de l'auteur}
tel que \uncertain{$B$}\note{on attendrait $C'$ : ce sont les arêtes de $C'$ que le morphisme contracte.} soit l'ens des arêtes \add{de $\underline{G}'$} qui se « contractent en un point ». On a
\[
\pi_0(G') \xrightarrow{\ \sim\ } \pi_0(G)
\]
[donc $G$ est $0$-connexe ssi $G'$ l'est]. \struck{Pour que} Soit, pour $\forall\, \alpha \in S$ (ens des sommets de $\underline{G}$) $G'_\alpha$ le sous-graphe de $G'$ fibre en $\alpha$ (qui est une composante connexe du s-graphe $G'_{!}$ de $G'$ défini par $S'$ et $C'$, \struck{\ill{}}) et $\underline{G}'_\alpha$ le MD-graphe déduit de $G'_\alpha$ en considérant l'application « genre » $S'_\alpha \to \mathbb{N}$ induite par $g' : S' \to \mathbb{N}$, et en prenant comme ens de pts marqués
\[
I'_\alpha \amalg \widetilde{B}'_\alpha
\]
(où $I'_\alpha = I' \,|\, S'_\alpha$, $\widetilde{B}' = \widetilde{A}' \smallsetminus \widetilde{C}'$ avec $\widetilde{C}' = \widetilde{A}' \,|\, C'$, $\widetilde{B}'_\alpha = \widetilde{B}' \cap \widetilde{A}'_\alpha$). Alors $\underline{G}'$ \struck{satisfait la condition} \add{est stable} (condition de régularité MD) ssi les $\underline{G}'(\alpha)$ le sont -- et dans ce cas $\underline{G}$ est aussi stable. Par suite, si $\underline{G}'$ est un MD-graphe, donc $\underline{G}$ et les $\underline{G}'(\alpha)$ le sont aussi.

\page{48}
\emph{Rem.} On a généralisé par la bande les définitions de morphisme de généralisation entre MD-graphes (ou m-graphes \uncertain{quelconques}, ni \uncertain{0-connexes} ni stables), pour avoir une plus jolie \uncertain{terminologie}. Les m-graphes déduits de $\underline{G}'$ par le procédé décrit s'appellent les \emph{m-graphes généralisations} de $\underline{G}'$, ou \emph{m-graphes quotients}. Elles forment un ens.\ ordonné, qui est isomorphe, par $\underline{G} \longleftrightarrow C'$, à l'opposé de l'ens ordonné des parties $C'$ de $A'$ (ens des arêtes de $\underline{G}'$), donc, par
\[
\underline{G} \longleftrightarrow B' = A' \smallsetminus C',
\]
à l'ens ordonné des parties $B'$ de $A'$. Donc si $\mu' = \operatorname{card} A'$, le nb de généralisations, de $\underline{G}'$ est $2^{\mu'}$.

Appelons « \emph{écart} » d'un morphisme de généralisation $\underline{G}' \to \underline{G}$ de m-graphes l'entier naturel $\delta = \operatorname{card} C'$. Alors l'écart \struck{est \ill{}} \add{(pour $\underline{G}'$ fixé)} \struck{est} peut prendre toutes les valeurs, entre $0$ et $s_1(\underline{G}') =$ codim modulaire de $\underline{G}'$. Il est nul ssi le morphisme de gén.\ est un \emph{iso}.

\page{49}\note{p. 24 de l'auteur}
Tout morphisme de généralisation se factorise en un composé de morphismes de généralisation d'écart~1. \struck{de façon unique} \add{Le nb de} façons de factoriser ainsi, à isom.\ près, est égal à $\delta!$, où $\delta$ est l'écart du morphisme de généralisation.

Considérons, pour un MD-graphe donné $\underline{G}$, de type \add{numérique} $(g, \nu)$\note{avant $\nu$, un symbole biffé illisible.}, l'\struck{les} \add{inclusion des} schémas modulaires propres sur $S$
\[
(57)\qquad \widehat{M}_{[\underline{G}]} \hookrightarrow \widehat{M}_{g,\nu}
\]
et aussi l'inclusion
\[
(58)\qquad (\widehat{M}_G, \Gamma^{!}) \hookrightarrow \widehat{M}_{g,I}
\]
\note{sous $(\widehat{M}_G, \Gamma^{!})$, une accolade et une notation biffée, avec $G$ en indice.}
où
\[
(59)\qquad
\left\{
\begin{array}{l}
\Gamma^{!} \subset \Gamma = \operatorname{Aut}(\underline{G}) \\[2pt]
\Gamma^{!} \overset{\text{déf}}{=} \{\, u \in \Gamma \mid u_I = \mathrm{id}_I \,\} = \operatorname{Ker}\bigl(\operatorname{Aut}(\underline{G}) \to \mathfrak{S}_I\bigr),
\end{array}
\right.
\]
où $I$ est l'ens.\ des « pts marqués » de $\underline{G}$.

Alors les \add{$2^{\mu}$} généralisations de $\underline{G}$ sont en correspondance biunivoque avec l'ens.\ des sous-multiplicités (fermées, \struck{\ill{}} donc propres sur $\mathbb{Z}$, et lisses sur $\mathbb{Z}$) de $\widehat{M}_{g,\nu}$ (resp.\ de $\widehat{M}_{g,I}$) de la forme $\widehat{M}_{[\underline{G}']}$ (resp.\ $(\widehat{M}_{\underline{G}'}, \Gamma'^{!})$) qui contiennent $\widehat{M}_{[\underline{G}]}$ (resp.\ $(\widehat{M}_G, \Gamma^{!})$). C'est

\page{50}
la situation typique d'une stratification : l'infini « à croisements normaux », et l'analyse locale au voisinage d'une \add{MD-}courbe géométrique $(\hat{X}, I)$ donnée (sur un corps alg.\ clos), à coups de « schémas modulaires \struck{locaux} formels » montre qu'on a bel et bien une telle situation à croisements normaux (on y reviendra par la suite).

\section{7. Étude des $\widehat{M}_G$ de dim $\leqslant 2$}

\page{51}\note{p. 25 de l'auteur. Le titre est de sa main, en tête de page : « Étude des \struck{\ill{}} $\widehat{M}_G$ \struck{\ill{}} \emph{de dim} $\leqslant 2$ ».}
a) Détermination des MD-graphes correspondants.

I) \emph{Cas de la dimension 0.} \struck{Les MD-gr} \add{conditions} suivantes sur le MD-graphe $\underline{G}$ sont équivalentes

a) $\dim M_G = 0$ \quad ($\Leftrightarrow \dim \widehat{M}_G = 0$)

b) \struck{\ill{}} La « dim.\ modulaire » de $\underline{G}$ est nulle, i.e.\ le nb $s_1$ d'arêtes de \struck{\ill{}} $\underline{G}$ (la codimension modulaire de $\underline{G}$) est égal à son maximum $3(g-1)+\nu$.

c) \struck{\ill{}} $\underline{G}$ est maximal parmi les MD-graphes pour la relation de généralisation, i.e.\ tout morphisme de généralisation $\underline{G}' \to \underline{G}$ de MD-graphes est un iso.

d) Tout sommet \add{$\alpha$} de $\underline{G}$ est de genre $0$, et est d'ordre \struck{\ill{}} $\omega_\alpha \leqslant 3$ \struck{i.e.\ le poids marqué} le \add{poids} $\nu_\alpha$ est égal à $3 - \omega_\alpha$.

On peut donc, \emph{à isomorphisme unique près}, \struck{trouver} une courbe standard \add{$\underline{G}$-épinglée}, et une seule sur une base quelconque $\mathcal{S}$, \struck{$G$-épinglée}, i.e.
\[
(60)\qquad M_G = M_{\widehat{G}} \simeq \operatorname{Spec} \mathbb{Z} \overset{\text{déf}}{=} S_0,
\]
donc
\[
(61)\qquad M_{[G]} \simeq B(\Gamma_{S_0})
\]
(\struck{\ill{}} multiplicité classifiante \add{pour le} schéma en groupes $\Gamma_{S_0}$ sur $S_0 = \operatorname{Spec} \mathbb{Z}$, où $\Gamma = \operatorname{Aut} \underline{G}$).

\page{52}
\[
(62)\qquad (M_{[G]}, \Gamma^{!}) \simeq B(\Gamma^{!}_{S_0})
\]
La « dimension modulaire \add{ambiante} » d'un MD $\underline{G}$ est égale à \struck{$\ldots \leqslant 2h_1(G) + \nu$}\note{formule biffée, en partie illisible ; sous le $\nu$ biffé : « $= \operatorname{card} I$ ».} :
\[
(63)\qquad \operatorname{dimmod}(G) = 3(g-1) + \nu = 3\bigl(h_1(G) - 1\bigr) + \nu .
\]
\note{sous $\nu$ : « $= \operatorname{card} I$ » ; en marge : « puisque $g = \sum_\alpha g_\alpha + h_1$, $g_\alpha = 0$ ».}

Nous allons faire l'énumération des cas où cette dimension est $\leqslant 2$, ce qui implique déjà $h_1 \leqslant 1$, i.e.\ $h_1 \in \{0, 1\}$.

\struck{a) $\operatorname{dimmod}(G) = 0$}

a) $h_1 = 0$, i.e.\ $G$ est un \emph{arbre}, la dim.\ modulaire est $\nu - 3$, dans ce cas $\nu \geqslant 3$. Si on veut que cette dim soit $\leqslant 2$, cela signifie que $3 \leqslant \nu \leqslant 5$.

Un arbre qui n'est pas réduit à un \struck{\ill{}} sommet, ni à un segment, a au moins trois bouts, dont chacun doit avoir un poids marqué $\geqslant 2$, donc $\nu \geqslant 6$ -- donc si $\nu \leqslant 5$, on voit que l'arbre est \uncertain{réduit} à un sommet, ou un segment. On trouve les cas
\[
(64)\qquad
\left\{
\begin{array}{lll}
\bullet^{3} & \nu = 3, & \dim \operatorname{amb} G = 0 \\[4pt]
\overset{2}{\bullet}\!\!-\!\!\!-\!\!\overset{2}{\bullet} & \nu = 4, & \dim \operatorname{amb} G = 1 \\[4pt]
\overset{2}{\bullet}\!\!-\!\!\!-\!\!\overset{1}{\bullet}\!\!-\!\!\!-\!\!\overset{2}{\bullet} & \nu = 5, & \dim \operatorname{amb} G = 2
\end{array}
\right.
\]
\drawing{En regard de chacun des trois cas de (64), la courbe correspondante : une droite portant trois points marqués $a, b, c$ ; deux droites sécantes portant $a, b$ et $c, d$ ; une chaîne de trois droites, les extrêmes portant $a, b$ et $d, e$, celle du milieu portant $c$.}

\page{53}\note{p. 26 de l'auteur}
b) $h_1 = 1$, donc $\operatorname{dimmod} G = \nu$, on veut $0 \leqslant \nu \leqslant 2$. On voit que le graphe $G$ a au maximum un bout, tous les autres sommets étant donc d'ordre~3. \struck{\ill{}} Mais un graphe (\add{top.\ connexe}) avec $h_1 = 1$ et n'ayant qu'un seul bout top.\ est homéomorphe à
\drawing{un cercle muni d'une queue : un segment attaché au cercle et terminé par un sommet libre.}
et s'il \struck{\ill{}} \add{n'a pas de} sommet d'ordre~2, c'est qu'il est \add{\uncertain{le}} graphe \struck{\ill{}} \add{combinatoire} correspondant à la figure précédente. D'autre part, si le graphe n'a pas de bout, c'est qu'il est homéomorphe au cercle. Tous ses sommets sont d'ordre~2, donc de poids marqué~1, le nb de ces sommets \struck{de} est \add{dans} $1 \ldots 2$, (il doit y en avoir au moins~1) on a pour graphe \drawing{un cercle avec un sommet} et \drawing{un cercle avec deux sommets}. En somme, on trouve exactement trois cas, résumés en
\[
(65)\qquad
\left\{
\begin{array}{ll}
\text{boucle à un sommet, de poids } 1 & \nu = 1 = \dim \operatorname{amb} G \\[2pt]
\text{deux sommets de poids } 1 \text{ joints par deux arêtes} & \nu = 2 = \dim \operatorname{amb} G \\[2pt]
\text{boucle munie d'une queue, le bout de poids } 2 & \nu = 2 = \dim \operatorname{amb} G
\end{array}
\right.
\]
\drawing{Les trois graphes de (65) sont dessinés : un cercle portant un sommet marqué « 1 » ; deux sommets marqués « 1 » reliés par deux arcs ; un cercle avec une queue dont le bout est marqué « 2 ». En regard, les courbes correspondantes : une courbe nodale (en huit) portant un point $a$ ; deux droites se coupant en deux points, portant $a$ et $b$ ; une courbe nodale à laquelle est attachée une droite portant $a$ et $b$.}

\page{54}
\marginal{MD-graphes maximaux (dim mod.\ nulle), de dim.\ modulaire ambiante $\leqslant 2$}
Nous allons aussi déterminer les groupes modulaires $\Gamma$ et $\Gamma^{!}$ correspondant à ces six cas :
\[
(65)\qquad
\left\{
\begin{array}{lll}
\bullet^{3} & \Gamma \simeq \mathfrak{S}_3, \ \Gamma^{!} = \{1\} & (\dim \operatorname{amb} 0) \\[4pt]
\overset{2}{\bullet}\!\!-\!\!\!-\!\!\overset{2}{\bullet} & \Gamma \simeq D_4, \ \Gamma^{!} = \{1\} & (\dim \operatorname{amb} 1) \\[4pt]
\overset{2}{\bullet}\!\!-\!\!\!-\!\!\overset{1}{\bullet}\!\!-\!\!\!-\!\!\overset{2}{\bullet} & \Gamma \simeq D_4, \ \Gamma^{!} = \{1\} & (\dim \operatorname{amb} 2)
\end{array}
\right.
\]
\note{il numérote ici (65) le tableau des arbres, numéroté (64) p. 52 ; le tableau des graphes à $h_1 = 1$, numéroté (65) p. 53, devient (66).}
\[
(66)\qquad
\left\{
\begin{array}{ll}
\text{boucle à un sommet de poids } 1 : & \Gamma = \Gamma^{!} \simeq \mathbb{Z}/2 \quad (\dim \operatorname{amb} 1) \\[2pt]
\text{deux sommets de poids } 1, \text{ deux arêtes} : & \Gamma \simeq D_2, \ \Gamma^{!} \simeq \mathbb{Z}/2 \quad (\dim \operatorname{amb} 2) \\[2pt]
\text{boucle avec queue, bout de poids } 2 : & \Gamma \simeq \mathbb{Z}/2 \times \mathbb{Z}/2, \ \Gamma^{!} \simeq \mathbb{Z}/2 \quad (\dim \operatorname{amb} 2)
\end{array}
\right.
\]
\note{les graphes de (66) sont dessinés ; on les décrit ici en mots, comme ceux de (65) p. 53.}

Notons que pour un \add{MD-}graphe $G$ de dim modulaire nulle, si $(\mathcal{X}, I_{\mathcal{S}})$ est une \add{MD-}courbe \struck{\ill{}} $G$-épinglée sur une base $\mathcal{S}$, alors le groupe des \uncertain{automorphismes} de $(\mathcal{X}, I_{\mathcal{S}})$ (\uncertain{qui ne} respectent pas forcément l'épinglage) est can.\ isom.\ à $\Gamma$, et \struck{celui des} le \struck{$\mathcal{X}$,} sous-groupe qui induit l'identité sur $I_{\mathcal{S}}$ est can.\ isom.\ à $\Gamma^{!}$.

\page{55}\note{p. 27 de l'auteur}
Revenons au cas d'un MD-graphe général $\underline{G}$ ; sa dimension modulaire est
\[
(67)\qquad \delta(G) = \sum_{\alpha \in S} \underbrace{\bigl(3(g_\alpha - 1) + \nu_\alpha + \tilde{\mu}_\alpha\bigr)}_{\delta_\alpha}
\]
si on veut qu'elle soit égale à $1$ ou $2$, cela signifie que l'on est dans l'un des deux cas :
\[
\left\{
\begin{array}{l}
1 \text{ ou } 2 \text{ des } \delta_\alpha \text{ sont égaux à } 1, \text{ les autres sont nuls} \\
1 \text{ des } \delta_\alpha \text{ est égal à } 2, \text{ les autres sont nuls}
\end{array}
\right.
\]
Le cas $\delta(G) = 1$ correspond au cas où exactement un des $\delta_\alpha$ est $= 1$, les autres étant nuls.

II) \emph{$\underline{G}$ de dim modulaire 1}

Cela signifie qu'il existe un sommet $\alpha_0$ (et un seul, \struck{de} fait) de $\underline{G}$ qui \add{satisfasse la} \struck{soit de genre} condition
\[
(68)\qquad
\left(
\begin{array}{l}
\alpha_0 \text{ de genre } 1, \text{ et son poids total } \nu_{\alpha_0} + \tilde{\mu}_{\alpha_0} \text{ est égal à } 1 \\
\quad \text{i.e.\ } \alpha_0 \text{ est d'ordre } 1, \text{ ou isolé de poids marqué } 1 \\
\underline{\text{ou}} \ \bigl(\alpha_0 \text{ de genre } 0, \text{ et son poids total } \nu_{\alpha_0} + \tilde{\mu}_{\alpha_0} = 4, \\
\quad \text{i.e.\ } (\nu_{\alpha_0}, \tilde{\mu}_{\alpha_0}) \in \{(4,0), (3,1), (2,2), (1,3), (0,4)\}\bigr)
\end{array}
\right)
\]
avec en plus
\[
(69)\qquad
\begin{array}{l}
\text{les sommets } \alpha \neq \alpha_0 \text{ sont de genre } 0, \text{ d'ordre } \tilde{\mu}_\alpha \leqslant 3, \\
\text{et pour un tel } \alpha \text{ le poids marqué } \nu_\alpha \text{ est égal à } 3 - \tilde{\mu}_\alpha .
\end{array}
\]
\struck{On voit que la MD-structure d'un tel graphe}

\page{56}
\marginal{MD-graphes élémentaires (sommet isolé) de dim.\ modulaire 1 \add{i.e.\ codim modulaire nulle} \struck{ambiante 1}}
\struck{est déterminée par la structure de graphe et des sommets prédéfinis par le}

Les cas les plus simples sont les cas élémentaires (graphes \struck{\ill{}} \add{ponctuels}) :
\[
(70)\qquad
\left\{
\begin{array}{ll}
\overset{4}{\bullet} \ \ g = 0 & \Gamma = \mathfrak{S}_4, \ \Gamma^{!} = 1 \\[6pt]
\overset{1}{\bullet} \ \ g = 1 & \Gamma = \Gamma^{!} = \{e\}
\end{array}
\right.
\]

Les autres cas sont ceux des \add{MD-}graphes « \emph{rétro-nuls} » (tous les $g_\alpha$ nuls) tels que tous les $\hat{\nu}_\alpha$ soient égaux à trois (donc les sommets d'ordre $1 \leqslant \tilde{\mu}_\alpha \leqslant 3$) \struck{et} sauf l'un d'eux, qui est égal à $4$ (donc $1 \leqslant \tilde{\mu}_{\alpha_0} \leqslant 4$), ou bien des \add{MD-}graphes qui \uncertain{soient} tous leurs sommets sauf un rétro-nuls et $\hat{\nu}_\alpha = 3$, \struck{donc} sauf un \add{$\alpha_0$} qui est un sommet-bout ($\tilde{\mu}_\alpha = 1$) et de genre $1$, $\nu_{\alpha_0} = 0$ -- donc ici tous les sommets sont d'ordre $1 \leqslant \tilde{\mu}_\alpha \leqslant 3$, et les $\nu_\alpha$ \add{pour $\alpha \neq \alpha_0$} sont égaux à $3 - \tilde{\mu}_\alpha$, \struck{pour $\alpha \neq$} et $\nu_{\alpha_0} = 0$.

Dans l'un et l'autre cas \add{($g_{\alpha_0} = 0$, $g_{\alpha_0} = 1$)} on aura \struck{un} iso can.\ de multiplicités
\[
(71)\qquad \widehat{M}_G \simeq \widehat{M}_{g_{\alpha_0}, \hat{I}_{\alpha_0}}
\]
avec $\hat{I}_{\alpha_0} = I_{\alpha_0} + \widetilde{A}_{\alpha_0}$ (de cardinal $4$ ou $1$ suivant que $g_{\alpha_0} = 0$ ou $g_{\alpha_0} = 1$).

\page{57}\note{p. 28 de l'auteur}
\marginal{MD-graphes de dim.\ modulaire 1 \uncertain{et} de dim.\ modulaire ambiante \uncertain{1, 2}}
\struck{\ill{}} et l'action de $\Gamma_{\underline{G}} = \Gamma$ sur $\widehat{M}_G \simeq \widehat{M}_{g_{\alpha_0}, \hat{I}_{\alpha_0}}$ est triviale si $g_{\alpha_0} = 1$, et elle se fait à travers
\[
(72)\qquad \Gamma_{\underline{G}} \longrightarrow \mathfrak{S}_{\hat{I}_{\alpha_0}} \quad (\simeq \mathfrak{S}_4)
\]
dans le cas où $g_{\alpha_0} = 0$. L'image de $\Gamma_{\underline{G}}$ dans $\mathfrak{S}_4$ doit pouvoir être n'importe quel sous-groupe de $\mathfrak{S}_4$, sauf erreur. Explicitons déjà que sont les cas où la dim.\ ambiante est $\leqslant 2$. Le cas de la dim ambiante $1$, i.e.\ où la codim modulaire \add{$s_1$} de $\underline{G}$ est $0$, est évidemment celui donné par (70), i.e.\ $G$ est réduit à un sommet isolé.

Considérons donc le cas de la dim mod.\ ambiante $2$, i.e.\ codim modulaire $1$, i.e.\ \emph{une} arête. On trouve soit \drawing{une boucle sur un sommet}, soit \drawing{un segment}.

On trouve les trois cas de dim modulaire \struck{1} ambiante $2$, dim modulaire $1$
\struck{Dans le cas $g_{\alpha_0} = 0$ tous les sommets sauf $\ldots$ sont d'ordre $\leqslant$}
\[
(72)\qquad
\left\{
\begin{array}{lll}
\overset{3}{\underset{\alpha_0}{\bullet}}\!\!-\!\!\!-\!\!\overset{2}{\underset{\alpha}{\bullet}} & \text{i.e.\ deux droites } (a, b, c;\ d, e) & \text{type } (0,5) \\[6pt]
\text{boucle sur } \alpha_0, \text{ de poids } 2 & \text{i.e.\ courbe nodale } (a, b) & \text{type } (1,2) \\[6pt]
\overset{}{\underset{(\alpha_0),\ g_{\alpha_0} = 1}{\bullet}}\!\!-\!\!\!-\!\!\overset{2}{\underset{g_\alpha = 0}{\bullet}} & \text{i.e.\ genre } 1 + \text{droite } (a, b) & \text{type } (1,2)
\end{array}
\right.
\]
\note{le numéro (72) est porté deux fois sur la page. Les graphes et les courbes sont dessinés ; les courbes sont décrites ici en mots. Dans le premier graphe, l'indice du second sommet est surchargé ($\alpha_0$ ?). Au-dessus du sommet $\alpha_0$ du troisième graphe, un chiffre biffé.}

(NB les sommets \add{à poids marqué} des branches des courbes, dont le genre \add{n'est pas marqué ou poids marqué} est marqué en dessous, sont de genre \struck{zéro} \add{\ill{}} ; les poids marqués sont indiqués sur les \struck{graphes} \add{au-dessus des} sommets)

\page{58}
Les groupes d'automorphismes de ces MD-graphes sont
\[
(73)\qquad
\left\{
\begin{array}{ll}
\overset{3}{\underset{\alpha_0}{\bullet}}\!\!-\!\!\!-\!\!\overset{2}{\bullet} & \Gamma = \mathfrak{S}_3 \times \mathfrak{S}_2, \ \Gamma^{!} = 1, \ \operatorname{Im}(\Gamma \to \mathfrak{S}_4) \simeq \mathfrak{S}_3 \\[6pt]
\text{boucle sur } (\alpha_0), \text{ de poids } 2 & \Gamma \simeq \mathfrak{S}_2 \times \mathfrak{S}_2, \ \Gamma^{!} \simeq \mathfrak{S}_2, \ \operatorname{Im}(\Gamma \to \mathfrak{S}_4) \simeq \Gamma \\ & \qquad = \mathfrak{S}_2 \times \mathfrak{S}_2 \\[6pt]
\underset{g_{\alpha_0} = 1}{\bullet}\!\!-\!\!\!-\!\!\overset{2}{\bullet} & \Gamma \simeq \mathfrak{S}_2, \ \Gamma^{!} \simeq 1
\end{array}
\right.
\]
\note{après $\Gamma^{!} \simeq 1$ à la dernière ligne, une suite de signes biffés.}

III) \emph{Dim.\ modulaire 2}

Correspond par (67), aux deux cas suivants :

1°) Tous les $\alpha \in S$ sauf un seul \add{$\alpha_0$} sont de genre $0$, et de poids total $3$ (i.e.\ d'ordre \add{$\tilde{\mu}_\alpha$} $\leqslant 3$, et de poids marqué $3 - \tilde{\mu}_\alpha$), et $(\alpha_0, \hat{\nu}_{\alpha_0})$ est d'un des deux types $(0,5)$, $(1,2)$ \struck{\ill{}}, i.e.\ $\alpha_0$ est de genre $0$, d'ordre \add{$\tilde{\mu}_{\alpha_0}$} $\leqslant 5$ et de poids marqué $5 - \tilde{\mu}_{\alpha_0}$, \emph{ou} $\alpha_0$ de genre $1$, d'ordre \add{$\tilde{\mu}_\alpha \in \{1, 2\}$} \struck{\ill{}} et de poids marqué $2 - \tilde{\mu}_{\alpha_0} \in \{1, 0\}$.
\marginal{donc deux cas en réalité}

2°) Tous les $\alpha \in S$ sauf deux exactement $\alpha_0$ et $\alpha_1$ sont de genre $0$, \add{et} de poids total $3$ (i.e.\ d'ordre $1 \leqslant \tilde{\mu}_\alpha \leqslant 3$, et de poids \struck{total} marqué égal à $3 - \tilde{\mu}_\alpha$), $\alpha_0$ et $\alpha_1$ sont \struck{\ill{}} chacun de type $(g_i, \hat{\nu}_i)$ égal soit à $(0,4)$, soit à $(1,1)$.
\marginal{Cela fait donc trois cas, suivant que les deux types \uncertain{sont} \ill{} $(0,4)$, $(1,1)$, \ill{} $(0,4)$, \ill{} $(1,1)$.}

La \add{multiplicité} \struck{schéma} modulaire $\widehat{M}_G$ est donc isomorphe, suivant les cas, à l'une des suivantes :

\page{59}\note{p. 29 de l'auteur}
\marginal{\struck{\ill{}} multiplicités modulaires de dim $2$ \ill{} $\widetilde{M}_G$}
\[
(74)\qquad
\left\{
\begin{array}{lll}
\widehat{M}^{!}_{0,5} & \Gamma \text{ opère via} & \Gamma \to \mathfrak{S}_{\hat{I}_{\alpha_0}} \ (\simeq \mathfrak{S}_5) \\[2pt]
\widehat{M}^{!}_{1,2} & \Gamma \text{ ------} & \Gamma \to \mathfrak{S}_{\hat{I}_{\alpha_0}} \simeq \mathbb{Z}/2\mathbb{Z} \\[2pt]
\widehat{M}^{!}_{0,4} \times \widehat{M}^{!}_{0,4} & \Gamma \text{ ------} & \Gamma \to \mathfrak{S}_{\hat{I}_{\alpha_0} \amalg \hat{I}_{\alpha_1}} \simeq \mathfrak{S}_8 \\[2pt]
\widehat{M}^{!}_{0,4} \times \widehat{M}^{!}_{1,1} & \Gamma \text{ ------} & \Gamma \to \mathfrak{S}_{\hat{I}_{\alpha_0}} \simeq \mathfrak{S}_4 \\[2pt]
\widehat{M}^{!}_{1,1} \times \widehat{M}^{!}_{1,1} & \Gamma \text{ opère trivialement.} &
\end{array}
\right.
\]
\note{les tirets sont ses « idem ». À la 3\textsuperscript{e} ligne, un symbole biffé avant $\mathfrak{S}$ ; à la 4\textsuperscript{e}, une surcharge (« $\alpha \in \ldots$ » ?) avant $\mathfrak{S}_{\hat{I}_{\alpha_0}}$.}

\emph{NB} Dans le cas 3°) de $\widehat{M}_{0,4} \times \widehat{M}_{0,4}$, $\Gamma$ \add{s'envoie} \struck{opère} dans $\mathfrak{S}_8$ par l'intermédiaire de
\[
\Gamma \longrightarrow \underbrace{\mathbb{Z}/2 \cdot (\mathfrak{S}_4 \times \mathfrak{S}_4)}_{\text{prod.\ }\tfrac{1}{2}\text{ direct}} .
\]

On peut considérer que pour l'essentiel, les \add{seules} \struck{\ill{}} multiplicités \add{$\widehat{M}_G$} modulaires de dim $1$, \struck{\ill{}} \uncertain{nouvelles} sont :
\[
(75)\qquad
\left\{
\begin{array}{ll}
M^{!}_{0,4} & \Gamma \text{ opère via } \Gamma \to \mathfrak{S}_{I_{\alpha_0}} \simeq \mathfrak{S}_4 \\[2pt]
\widehat{M}^{!}_{1,1} & \Gamma \text{ opère trivialement,}
\end{array}
\right.
\]
\marginal{(cas élémentaires)}
les « cas nouveaux » dans (74) sont $\widehat{M}^{!}_{0,5}$ et $\widehat{M}^{!}_{1,2}$. Cela fait en tout essentiellement quatre cas « élémentaires » à étudier à fond, pour le cas de dim.\ modulaires $\leqslant 2$. J'ai l'impression qu'une compréhension plus ou moins complète de la structure des $\Pi_1$ discrets (i.e.\ dans le contexte analytique) impliqués dans ces cas, donnera la clef pour une compréhension complète dans le cas général.

\page{60}
Partant des MD-graphes « élémentaires » correspondants, je voudrais expliciter dans les quatre cas la structure \add{de spécialisation entre} des strates des $\widehat{M}_{g,\nu}$ et $\widehat{M}^{!}_{g,\nu}$.

\section{8) Structure de spécialisation des strates en dim modulaire $\leqslant 2$}

1°) \emph{Cas des $\widehat{M}^{!}_{0,4}$ et $\widehat{M}_{0,4}$}

Commençons par $\widehat{M}^{!}_{0,4}$.
\[
(76)\qquad
\begin{array}{c|c}
\text{Codim } 0 & \text{Codim } 1 \\ \hline
\overset{I(4)}{\bullet} \quad \xrightarrow{\ 3\ } & \overset{I'(2)}{\bullet}\!\!-\!\!\!-\!\!\overset{I''(2)}{\bullet}
\end{array}
\]
Le diagramme signifie que l'ens.\ marqué $I$ de cardinal $4$, \add{se partageant,} suivant l'une quelconque des trois partitions de $I$ en $I'$ et $I''$ de cardinal $2$ chacune, donne $3$ spécialisations différentes de la strate maximale -- cela correspond bien au \uncertain{fait} \uncertain{qu'on a} trois pts à l'infini de $U_{0,3}$, \uncertain{où} on a
\[
(77)\qquad \widehat{M}^{!}_{0,4} \simeq \mathbb{P}^1_{\mathbb{Z}}, \quad M^{!}_{0,4} \simeq (U_{0,3})_{\mathbb{Z}} = \mathbb{P}^1_{\mathbb{Z}} \smallsetminus \{0, 1, \infty\}
\]
Pour bien faire, il faudrait aussi expliciter
\note{la page s'arrête ici ; la phrase se poursuit au-delà de la fin du lot.}

\end{document}
